\paragraph{window-\texorpdfstring{$6$}{6} 残基半环与奇偶角色空间的计数分离}
在 \(m=6\) 的具体窗口上，折叠同余商的有限残基算术与 fold-gauge 的奇偶角色空间可以直接并列比较，从而得到一条严格的有限状态分离结论。

\begin{corollary}[window-\texorpdfstring{$6$}{6} 残基半环无法忠实承载全部奇偶角色]\label{cor:xi-window6-residue-semiring-parity-character-gap}
记
\[
\mathcal R_6:=\Omega_6/\!\equiv_6.
\]
则
\[
\mathcal R_6\cong \ZZ/(F_8\ZZ)=\ZZ/(21\ZZ),
\qquad
|\mathcal R_6^\times|=\varphi(21)=12,
\]
\[
|\operatorname{Idem}(\mathcal R_6)|=2^{\omega(21)}=4,
\qquad
\#\{x\in \mathcal R_6:\ x\neq 0,\ x \text{ 为幂零元}\}=0.
\]
并且全部幂等元恰为
\[
\operatorname{Idem}(\mathcal R_6)=\{\overline 0,\overline 1,\overline 7,\overline{15}\}.
\]

另一方面，
\[
\bigl(G_6^{\mathrm{bin}}\bigr)^{\mathrm{ab}}\cong (\ZZ/2\ZZ)^{21},
\qquad
\bigl|\Hom(G_6^{\mathrm{bin}},\{\pm 1\})\bigr|=2^{21}.
\]
因此，任何仅通过残基半环 \(\mathcal R_6\) 取值的读出都不可能分离全部奇偶角色；特别地，不存在单射
\[
\Hom(G_6^{\mathrm{bin}},\{\pm 1\})\hookrightarrow \mathcal R_6.
\]
\end{corollary}
\begin{proof}
由定理 \ref{thm:xi-fold-congruence-semiring-singularity-nilpotent-profile}，
\[
\mathcal R_6\cong \ZZ/(F_8\ZZ)=\ZZ/(21\ZZ),
\]
并且由于 \(21=3\cdot 7\) 为平方自由数，有
\[
|\mathcal R_6^\times|=\varphi(21)=12,
\qquad
|\operatorname{Idem}(\mathcal R_6)|=2^{\omega(21)}=4,
\]
且不存在非零幂零元。

为求幂等元的显式代表，只需利用中国剩余同构
\[
\ZZ/(21\ZZ)\cong \ZZ/(3\ZZ)\times \ZZ/(7\ZZ).
\]
幂等元在每个局部分量上只能取 \(0\) 或 \(1\)，故四个可能的 CRT 坐标为
\[
(0,0),\qquad (1,1),\qquad (1,0),\qquad (0,1).
\]
把它们回拉到 \(\ZZ/(21\ZZ)\) 中，恰得到
\[
\overline 0,\qquad \overline 1,\qquad \overline 7,\qquad \overline{15}.
\]

另一方面，由定理 \ref{thm:xi-foldbin-gauge-representation-distribution-law-m6-spectrum}，
\[
\bigl(G_6^{\mathrm{bin}}\bigr)^{\mathrm{ab}}\cong (\ZZ/2\ZZ)^{21},
\qquad
\bigl|\Hom(G_6^{\mathrm{bin}},\{\pm 1\})\bigr|=2^{21}.
\]
而
\[
|\mathcal R_6|=21<2^{21}.
\]
因此，从奇偶角色空间到 \(\mathcal R_6\) 的任何映射都不可能是单射。于是，任何仅通过 \(\mathcal R_6\) 取值的残基读出都无法忠实区分全部奇偶角色。证毕。
\end{proof}

\endinput
